\documentclass[11pt]{article}
\usepackage{fontspec}
\usepackage[french,english]{babel}
\usepackage{amsmath}
\usepackage{amssymb}
\usepackage{amsthm}
\usepackage{mathrsfs}
\usepackage{tikz-cd}
\usepackage{enumitem}

\newcommand{\printedpage}[1]{\marginpar{\footnotesize p.~#1}}
\newcommand{\msp}[1]{\marginpar{\footnotesize [#1]}}

\newcommand{\Fib}{\mathrm{Fib}}
\newcommand{\Fibf}{\mathrm{Fibf}}
\newcommand{\Parf}{\mathrm{Parf}}
\newcommand{\Catfib}{\mathrm{Catfib}}
\newcommand{\Fibsc}{\mathrm{Fibsc}}
\newcommand{\Cat}{\mathrm{Cat}}
\newcommand{\HOT}{\mathrm{HOT}}

\begin{document}

\printedpage{19}
% head: A. Grothendieck : Dérivateurs, Ch. XVI, Édition des Universités de Montpellier II et Paris VII

\section*{4\quad Axiom Loc (7) of $W$-fibrations}

4.1. Recall the notion of a \textit{weak $W$-fibration} (XII p. 51, 52): the strongest version of the 8 equivalent conditions for $f : X \longrightarrow S$ to be a weak fibration is this.

\begin{itemize}[leftmargin=2cm]
\item[\textbf{(Ff)}] $\forall\ S'' \longrightarrow S'$ in $\Cat/S$, with $S'$, $S''$ \textit{contractible,} the corresponding morphism $X'' = X \times_S S'' \longrightarrow X' = X \times_S S'$ is in $W$.
\end{itemize}

When $W$ is strongly saturated, then this implies the condition, apparently stronger:

\begin{itemize}[leftmargin=2cm]
\item[\textbf{(Ff$'$)}] $\forall\ S'' \longrightarrow S'$ in $\Cat/S$, such that $S'' \longrightarrow S'$ is a homotopism, $X'' \longrightarrow X'$ is in $W$.
\end{itemize}

I denote by $\Fibf_W$ the set of weak fibrations. I summarize in a diagram of inclusions the main classes of arrows in Cat that we have encountered, akin to classes of ``fibrations'' (for $W$, in a suitable sense):

\msp{page 28}

\[
\begin{tikzcd}
 & & \Fib_W \arrow[d] & & \\
\Parf_\omega \arrow[r] & \Parf_W \arrow[r] & \Fibf_W & \underbrace{\Catfib^W}_{\mathcal{F}_W} \arrow[l] & \underbrace{\Catfib^\omega}_{\mathcal{F}_\omega} \arrow[l] \ ,
\end{tikzcd}
\]
where $\Fib_W$ denotes the $W$-fibrations ($f : S' \longrightarrow S$ is in $\Fib_W$ if and only if the base change functor $X \longmapsto X' = X \times_S S'$ from $\Cat/S$ to $\Cat/S'$ is compatible with the localizers $W_S$ and $W_{S'}$), Catfib the Cat-fibrations, $\Catfib^W$ (resp. $\Catfib^\omega$) the Cat-fibrations such that the base change functors are in $W$ (resp. are in $W_\omega$, \textit{cf.} def. below) \footnote[20]{N.B. $\Catfib_0^W = \Catfib \cap \Fibf_W$, $\Catfib_0^\omega = \Catfib \cap \Fibf_{W_\omega}$}, $\Parf_W$ the $W$-perfect morphisms, \textit{i.e.} $W$-smooth and $W$-proper, finally $\Parf_\omega$ the $W_\omega$-perfect morphisms, where $W_\omega$ is the smallest localizer satisfying Loc (1, 2, 3). Among the kinds of ``fibrations'' considered, it is only this last one, together with $\Catfib^\omega$, that is independent of $W$.

As such, it seems to me to play a role similar, dual, to that of Dwyer morphisms, and I presume that one can make them play a similar role. In any \msp{page 29} case, we know that every arrow $f$ of Cat factors as $pi$, with $i \in W_\omega$ and $p \in \Parf_\omega$, hence $i \in W$ for every $W$ satisfying Loc (1, 2, 3). We also know that $\Catfib^\omega$ satisfies the strong factorization statement, $f$ factors as $f = pi$, $i \in W_\omega$, $p \in \Catfib_0^\omega$. \footnote[21]{Thus the factorization property holds for all these classes, except perhaps $\Fib_W$.}

Let us also note that $\Parf_\omega$ and $\Parf_W$, as well as $\Fib_W$, are stable not only under base change (this is the case for the 5 classes considered, which moreover contain all isomorphisms), but also under composition. (Moreover, $\Fib_W$ is stable under direct factors, and I presume that the same holds for $\mathrm{Liss}_W$ and for $\mathrm{Prop}_W$, hence for $\Parf_W$, in particular for $\Parf_{W_\omega}$, but I have not taken the trouble to check it.) On the other hand, a priori $\Fibf_W$ is not in general stable under composition, and $\Catfib^W$, $\Catfib^\omega$ have no reason
% foot: Version 18 mars 2019
\printedpage{20}
% head: A. Grothendieck : Dérivateurs, Ch. XVI, Édition des Universités de Montpellier II et Paris VII
either to be stable under composition \footnote[22]{this should be made explicit by examples.}.

4.2. Finally, it seems more ``illuminating'' to introduce, instead of the well-determined basic localizer $W_\omega$, a basic localizer % [sic] déterminée -> déterminé, accord avec « localiseur » (masculin)
\[
W_0 \subset \mathrm{Fl}(\Cat),
\]
which will play the role of a sort of ``reference localizer'' for formulating certain properties of $W$. We will always assume that
\[
W \supset W_0.
\]

\msp{page 30}

Depending on the case, $W_0$ may be $W_\omega$, or some other basic localizer defined as the ``smallest'' satisfying given stability axioms. Or it may be $\mathbf{W}_\infty$ (the usual weak equivalences), when one is interested in localizers that contain $\mathbf{W}_\infty$ (the only ones, apparently, with which one has worked up to now in homotopy). It may also be $W$, and the statements obtained by specializing the results obtained to the case where $W = W_0$ are by no means tautological. The diagram on page 28 is then replaced by the following diagram of inclusions, more illuminating for our purpose:
\[
\begin{tikzcd}
 & \Fib_W \arrow[d] & \\
\Parf_W \arrow[r] & \Fibf_W & \mathcal{F}_W \arrow[l] \\
\Parf_{W_0} \arrow[r] \arrow[u] & \Fibf_{W_0} \arrow[u] & \mathcal{F}_{W_0} \arrow[l] \arrow[u] \\
 & \Fib_{W_0}, \arrow[u] &
\end{tikzcd}
\tag*{(4.2.1)}
\]
where for every basic localizer $W$, I put for short
\[
\mathcal{F}_W = \Catfib^W,
\]
which thus also defines $\mathcal{F}_{W_0}$. The class $\Fib_W$ \msp{page 31} should be regarded as the focus of our attention; most of the properties of $W$ that we are going to review here, will be expressed by an inclusion of the form
\[
\mathcal{F} \subset \Fib_W,
\]
where $\mathcal{F}$ is one of the seven other classes of diagram (4.2.1); with the exception at most of the case $\mathcal{F} = \Fib_{W_0}$, \textit{i.e.} of the inclusion $\Fib_{W_0} \subset \Fib_W$. As a property of $W$, this inclusion seems very weak and for that reason not of much use, except of course in the case where $\Fib_{W_0} \longrightarrow \Fibf_{W_0}$ is an equality, so that the inclusion considered is equivalent to $\Fibf_{W_0} \subset \Fib_W$ (and implies that $\Parf_{W_0}$ and $\mathcal{F}_{W_0}$ are also contained in $\Fib_W$). This on the other hand, as we shall see, is a fairly strong and very useful property of $W$, implying in particular the fact that $(\Cat, W, \Fib_W)$ is a category with fibrations.
% foot: Version 18 mars 2019

\printedpage{21}
% head: A. Grothendieck : Dérivateurs, Ch. XVI, Édition des Universités de Montpellier II et Paris VII

\msp{page 32}

For purposes of later reference, I am going to formulate now the various main aspects of \textit{the $W$-fibration axiom} Loc (7). The formulation which I had arrived at (\textit{cf.} I, VI, XII) used the notion which I now call ``\textit{weak} $W$-fibrations'' (the former $W$-fibration, now dethroned!), and in its strongest form was written as the inclusion % [sic] W-fibrations faible -> W-fibration faible, accord du nombre
\[
\Fibf_W \subset \Fib_W.
\]
This notion of weak $W$-fibration, just like that of $W$-perfect morphism, however useful they may be, nevertheless seems to me a little technical, and I prefer to give a formulation that calls only on more familiar notions, more elementary in appearance. This is why I am going to replace $\Fibf_W$ by the subclass $\mathcal{F}_W = \Catfib^W$, which (we shall see) in fact comes to the same thing.

\msp{page 33}

This class $\mathcal{F}_W$ is extremely natural, and even imposes itself in the yoga of derivators. Indeed, in the construction of $\HOT_W(S)$ for $S$ in Cat, one passes to a category of fractions of the category
\[
\Cat(S) = \underline{\mathrm{Hom}}(S, \Cat),
\]
by a localizer $W(S)$ (which we had generally denoted by $W_S^{\mathrm{g}}$ or $W_S^{\mathrm{u}}$ \ldots). This category $\Cat(S)$ is on the other hand isomorphic to the category
\[
\Fibsc\, S
\]
of fibered and split categories over $S$, with as arrows the $S$-functors which respect the splitting (and in particular cartesian) (\textit{cf.} VI, and especially the diagram p. 175). We have seen (\textit{loc. cit.}) that the localized category $\HOT_W(S)$ is equivalent to the localized category of the very closely related category
\[
\Fib\, S \qquad \text{(which I would now prefer to denote by $\Catfib\, S$ or $\Catfib_S$),}
\]
formed by the \textit{fibered} categories over $S$, with as \msp{page 34} arrows only the cartesian functors. (One should note however that the canonical functor $\Fibsc\, S \longrightarrow \Fib\, S$ is faithful but not fully, nevertheless it induces an equivalence on the localized categories, defining in the obvious way, by $W$-equivalence fiber by fiber, \textit{i.e.} as $W_S^{\mathrm{u}}$, the relevant localizer on $\Fib\, S$.) Thus for this reason of principle, and moreover throughout the development of the theory of $\HOT_W$, the category $\Fib\, S$ and thereby, the fibrant functors, have played an essential role. The class of arrows which we denote here by $\mathcal{F}_W$ or $\Catfib^W$ (one could also denote it by $\Fib^W$, were it not for the risk of confusion with $\Fib_W$) is then interpreted very naturally starting from there. An arrow $f : X \longrightarrow S$ is in $\mathcal{F}_W$ if and only if it is a) fibrant, and thereby defines in a particularly direct way an object of $\HOT_W(S)$, since up to $S$-equivalence, $X/S$ comes from an object of $\Fibsc\, S$ or again of $\Cat(S)$, \msp{page 35} and b) the object in $\HOT_W(S)$ defined by $f$ is in $\HOT_W^{\mathrm{lc}}(S)$, \textit{i.e.} in the essential image of the strictly full subcategory $\Fibsc_0^W S$ of $\Fibsc\, S$ (under the canonical functor
% foot: Version 18 mars 2019
\printedpage{22}
% head: A. Grothendieck : Dérivateurs, Ch. XVI, Édition des Universités de Montpellier II et Paris VII
$\Fibsc \longrightarrow \HOT_W(S)$), or what comes to the same thing (at least if $W$ is strongly saturated) that the functor % [sic] saturée -> saturé, accord avec W (« W fortement saturé » ailleurs)
\[
\begin{aligned}
S^{\circ} \quad &\longrightarrow \quad \mathrm{Hot}_W \\
s \quad &\longmapsto \quad \mathrm{hot}_W(X_s)
\end{aligned}
\]
be a ``locally constant'' functor. Thus, we shall regard the $f \in \mathcal{F}_W$ as the most economical and most natural arrows in Cat, which define locally constant relative $W$-homotopy types.

Let us also note that all the classes of arrows (4.2.1) are stable under base change (and they contain the isomorphisms). The categories $\Fib_W$, $\Parf_W$ and their $W_0$ variants are also stable under composition, on the other hand, this is not the case for $\Fibf_W$ and $\mathcal{F}_W$ and their $W_0$ variants, except under special axioms on $W$ or on $W_0$). It follows at once \msp{page 36} that if $\mathcal{F}$ is one of these classes, the condition % [sic] « sur W_0). » -> parenthèse fermante sans parenthèse ouvrante correspondante
\[
\mathcal{F} \subset \Fib_W
\]
is equivalent to this one:

\begin{itemize}[label={},leftmargin=3.5cm]
\item For every cartesian diagram
\[
\begin{tikzcd}
X \arrow[d, "f"'] & X' \arrow[l, "q"'] \arrow[d, "f'"] \\
S & S', \arrow[l, "p"]
\end{tikzcd}
\tag*{(4.2.2)}
\]
with $f \in \mathcal{F}$ and $p \in W$, we have $q \in W$.

(N.B. We also have $f' \in \mathcal{F}$, by the stability recalled above.)
\end{itemize}

4.3. Here then (\textit{cf.} below) is the relevant axiom, Loc (7) with its three aspects a, b, c \footnote[23]{I finally added a fourth F to it (for ``fibration'').}. We have seen (and shall see again) that we have the implications
\[
\text{a} \qquad \Longrightarrow \qquad \text{b} \qquad \Longrightarrow \qquad \text{c} \qquad (\Longrightarrow \qquad \text{F}).
\]
When we refer to Loc (7) without further specification, it will be understood that this is the strongest form Loc (7 a), which thus implies the other two. No less is needed to be completely at ease. But I recall that the weakened variant Loc (7b) is sufficient to develop the theory of the functors $f_!$ for the categories $\HOT_W(S)$ \footnote[24]{Here my memory is playing a trick on me - for this one has no need at all of the already very strong axiom Loc (7b), but only of Loc (3 bis), which does seem to be satisfied for all ``reasonable'' basic localizers. % [sic] qu'il soit vérifiée -> qu'il soit vérifié, accord avec « il »
} (\textit{cf.} VI).

\msp{page 37}

This is moreover a surprising thing, since via the general yoga of derivators associated with categories of models, one would rather expect that to have left adjoints $f_!$ of the inverse image functors $f^*$, one must make hypotheses on $W$ of type \textit{$W$-cofibration}
% foot: Version 18 mars 2019
\printedpage{23}
% head: A. Grothendieck : Dérivateurs, Ch. XVI, Édition des Universités de Montpellier II et Paris VII
(essentially, Loc (6 a)). It seems unlikely to me that Loc (7 b), or even Loc (7), implies Loc (6 a), or even its weaker variant Loc (6$'$ a) (axiom of the cocartesian square, or of Mayer-Vietoris, weak form). If on the other hand these axioms are not \textit{satisfied,} whereas Loc (7) is, there is a chance that this will give an example where there is a formalism of the $f_!$, without the left exactness axiom being satisfied.

As for the form Loc (7 c), weaker still, it is nevertheless sufficient (as we shall see below) to ensure that Cat has enough $W$-fibrations \msp{page 38} (a result which is by no means obvious, and could only become apparent to me after I had learned of the substantial lemma of K. S. \textsc{Brown}, in his beautiful work where he introduces categories with fibrations and with cofibrations).

We will also give below some very simple cases, where the condition Loc (7 b) already implies Loc (7) ($W$ strongly saturated, \textit{i.e.} Loc (1 bis)), or even where Loc (7 c) implies Loc (7) (the case where $W$ satisfies the direct limit axiom Loc (8)). I point out in passing that I know of no basic localizer satisfying the axiom Loc (8) (which already implies $W = W^{\natural}$, \textit{i.e.} Loc (1 bis)), without $W$ being strongly saturated (\textit{i.e.} even satisfying Loc (1 ter)).

\begin{itemize}[label={},leftmargin=0.6cm]
\item \textbf{Loc (7)}
\begin{itemize}[leftmargin=1.2cm]
\item[a)] We have $\mathcal{F}_W \subset \Fib_W$, \textit{i.e.} for every cartesian diagram (4.2.2) in Cat, with $f \in \mathcal{F}_W$ (\textit{i.e.} $f$ Cat-fibrant and weakly $W$-fibrant) and $p \in W$, we have $q \in W$.
\item[b)] For every commutative diagram
\[
\begin{tikzcd}[column sep=small]
X \arrow[rr, "f"] \arrow[dr, "p"'] & & Y \arrow[dl, "q"] \\
 & S &
\end{tikzcd}
\]
\msp{page 39}
with $p, q \in \mathcal{F}_W$ and $f \in W$, we have $f \in W_S^{\mathrm{u}}$ (or, what comes to the same thing since $X$, $Y$ are $W$-smooth over $S$, $f_s \in W$ for every $s \in \mathrm{Ob}\, S$) \footnote[25]{\textbf{N.B.} The proof of $B(\mathcal{F}_W) \Longrightarrow B(\Fibf_W)$ in fact shows that it suffices to postulate b) when $f$ is cartesian, and even when $X$, $Y$ are split over $S$ and $f$ compatible with the splitting.}.
\item[c)] We have
\[
\mathcal{F}_W \cap W \subset W^{\mathrm{u}},
\]
in other words, for every arrow $f : X \longrightarrow S$ in Cat such that $f \in \mathcal{F}_W$ (one could call these arrows Cat-$W$-fibrations, it being understood that they are also Cat-fibrations, but not necessarily $W$-fibrations), if $f \in W$, then $f \in W_S^{\mathrm{u}}$ (or, what comes to the same thing, the fibers $X_s$ of $f$ ($s \in \mathrm{Ob}\, S$) are $W$-aspheric).
\end{itemize}
\end{itemize}
% foot: Version 18 mars 2019

\end{document}
